

\section{Цель работы}

Целью данной лабораторной работы является освоение метода динамического
программирования для решения задачи оптимального раскроя материала, а также
приобретение навыков перебора допустимых схем раскроя (карт) и синтеза
оптимального плана выпуска изделий.

В рамках работы требуется:
\begin{itemize}
    \item освоить методику генерации всех максимальных карт раскроя листа
          заданного размера фигурами нескольких типов с учётом допустимых
          поворотов;
    \item понять, как отсеиваются доминируемые схемы раскроя, которые
          заведомо не могут улучшить решение;
    \item усвоить постановку задачи целочисленного линейного программирования
          для минимизации числа использованных листов при выполнении
          заданного плана выпуска;
    \item приобрести практические навыки решения этой задачи методом
          динамического программирования по вектору выпуска;
    \item научиться интерпретировать результат: какие карты и в каком
          количестве нужно использовать, чтобы получить минимальное число
          листов.
\end{itemize}

В конечном счёте выполнение работы должно сформировать практические навыки
применения методов дискретной оптимизации к реальным задачам
производственного планирования.

\clearpage

\clearpage




\clearpage

\section{Теоретическая часть}

Задача раскроя (cutting stock problem) — классическая задача дискретной
оптимизации, впервые поставленная Л.~В.~Канторовичем и позднее
формализованная Гилмором и Гомори как задача целочисленного линейного
программирования. Она возникает всюду, где из стандартных заготовок
требуется выкроить набор деталей с минимальными отходами или
минимальным числом использованных заготовок.

В данной работе рассматривается прямоугольный лист размером $R \times C$
клеток ($R = 3$, $C = 5$, всего 15 клеток). Имеется $n$ типов деталей;
для каждого типа $f$ заданы размеры $w_f \times h_f$, площадь
$s_f = w_f h_f$ и потребность $b_f$. Допускается поворот детали на
$90^\circ$, то есть деталь может занимать на листе прямоугольник как
$w_f \times h_f$, так и $h_f \times w_f$. Детали не должны пересекаться
и выходить за границы листа; разрезы производятся по границам клеток,
гильотинность не предполагается.

Требуется выбрать такие способы раскроя листов и такое их количество,
чтобы суммарный выпуск деталей каждого типа был не меньше $b_f$, а
общее число использованных листов было минимальным.


\textbf{Схемой раскроя} (картой) называется размещение деталей на одном
листе, при котором:
\begin{enumerate}
    \item все детали целиком лежат в пределах листа;
    \item детали попарно не пересекаются;
    \item размещение \emph{максимально}: ни одну деталь ни одного типа
          (ни в одном из двух поворотов) нельзя добавить, не сдвинув
          уже размещённые.
\end{enumerate}

Условие максимальности важно: не максимальная карта доминируется более
полной, поэтому в оптимальном решении такие карты не встречаются, и их
можно исключить из рассмотрения.

Формально состояние листа описывается функцией
$g : \{1,\ldots,R\} \times \{1,\ldots,C\} \to \{-1,0,\ldots,n\}$,
где $g(i,j) = -1$ означает пустую клетку, а $g(i,j) = f$ — клетку,
занятую деталью типа $f$. Карта максимальна тогда и только тогда,
когда ни один прямоугольник $w_f \times h_f$ или $h_f \times w_f$ не
помещается целиком в свободные клетки.

Каждая карта $j$ описывается вектором
$a_j = (a_{j1}, \ldots, a_{jn}) \in \mathbb{Z}_{\ge 0}^n$, где
$a_{jf}$ — количество деталей типа $f$ на этой карте. Занятая площадь
равна $S_j = \sum_f a_{jf} s_f \le R \cdot C$; величина
$R \cdot C - S_j$ — отходы карты.


Число размещений конечно, но быстро растёт. Для полного перебора
применяется рекурсивный поиск в глубину с мемоизацией состояний листа.
Состояние кодируется 64-битным целым: каждой клетке отводится 4 бита
(значения $0,\ldots,n$ после сдвига на $+1$), всего $4 \cdot R \cdot C$
бит (для листа $3 \times 5$ — 60 бит):
\[
\mathrm{key}(g) = \sum_{i=1}^{R} \sum_{j=1}^{C}
\bigl(g(i,j) + 1\bigr) \cdot 2^{\,4((i-1)C + (j-1))}.
\]
Множество посещённых состояний хранится в хеш-таблице, что исключает
повторный обход одной и той же конфигурации.

Алгоритм:
\begin{enumerate}
    \item если состояние уже посещалось — выйти;
    \item пометить состояние посещённым;
    \item если на листе есть хотя бы одна деталь и размещение
          максимально — сохранить карту в словаре с ключом $a_j$;
    \item для каждой детали и каждого поворота перебрать все позиции,
          в которых она помещается: поставить, рекурсивно продолжить,
          снять.
\end{enumerate}

Если несколько карт имеют одинаковый вектор $a_j$, но разное
расположение деталей, для оптимизации они эквивалентны — в модели
важны только количества. Поэтому из каждой группы сохраняется одна
карта-представитель.


В режиме <<не меньше>> карта $i$ \emph{доминируется} картой $j$, если
\[
a_{jf} \ge a_{if} \ \ \forall f, \qquad \exists f_0:\ a_{j f_0} > a_{i f_0}.
\]
Тогда любое использование $i$ можно заменить на $j$ без увеличения
числа листов и без уменьшения выпуска, поэтому карту $i$ можно
исключить. Геометрически это выделение Парето-оптимального фронта в
$n$-мерном пространстве векторов $a_j$. В режиме <<ровно>>
доминирование не применяется, так как перевыполнение там запрещено.

\clearpage

Введём переменные $x_j \in \mathbb{Z}_{\ge 0}$ — число листов,
раскроенных по карте $j$. Задача принимает вид
\[
\min \sum_{j=1}^{m} x_j
\]
при ограничениях
\[
\sum_{j=1}^{m} a_{jf} \, x_j \ge b_f, \quad f = 1, \ldots, n,
\qquad x_j \ge 0,\ x_j \in \mathbb{Z}.
\]
Это задача целочисленного линейного программирования с большим числом
переменных и малым числом ограничений. В классической постановке
Гилмора--Гомори она решается методом генерации столбцов; однако при
малом размере листа и небольших потребностях возможен более простой и
точный подход — динамическое программирование по вектору выпуска.


Пространство состояний имеет размер
$(b_A+1)(b_B+1)(b_C+1)(b_D+1) = 22 \cdot 16 \cdot 23 \cdot 32 = 258\,048$,
что позволяет решить задачу точно.

Обозначим через $dp[a][b][c][d]$ минимальное число листов, дающее
выпуск $a, b, c, d$ по типам A, B, C, D соответственно. В режиме
<<не меньше>> выпуск сверх потребности неотличим от выпуска ровно по
потребности, поэтому координаты ограничиваются значениями $b_f$:
\[
a \in \{0,\ldots,b_A\},\quad b \in \{0,\ldots,b_B\},\quad
c \in \{0,\ldots,b_C\},\quad d \in \{0,\ldots,b_D\}.
\]

Начальное состояние: $dp[0][0][0][0] = 0$, остальные — бесконечность.
Переход для карты $j$ с вектором $a_j = (a_{jA}, a_{jB}, a_{jC}, a_{jD})$:
\[
a' = \min(b_A,\ a + a_{jA}), \quad b' = \min(b_B,\ b + a_{jB}),
\]
\[
c' = \min(b_C,\ c + a_{jC}), \quad d' = \min(b_D,\ d + a_{jD}),
\]
\[
dp[a'][b'][c'][d'] = \min\bigl(dp[a'][b'][c'][d'],\ dp[a][b][c][d] + 1\bigr).
\]
Операция $\min(b_f, \cdot)$ реализует насыщение координаты. В режиме
<<ровно>> насыщение не применяется, а переходы, выводящие координату
за $b_f$, отбрасываются. Ответ задачи: $dp[b_A][b_B][b_C][b_D]$.

Восстановление плана выполняется по массивам предшественников: для
каждого состояния запоминаются номер карты и предыдущее состояние.
Пройдя от конечного состояния к начальному, получаем набор карт и их
кратностей $x_1, \ldots, x_m$.


Временная сложность алгоритма составляет
$O\bigl(m \cdot \prod_{f=1}^{n}(b_f+1)\bigr)$, где $m$ — число карт
после отсеивания доминируемых. В данной задаче $m = 21$, объём
пространства состояний $258\,048$, что даёт порядка $5 \cdot 10^6$
операций — доли секунды. Память — несколько массивов длины $258\,048$.

Корректность метода опирается на три свойства:
\begin{enumerate}
    \item \textbf{Принцип оптимальности Беллмана}: если оптимальное
          решение использует карту $j$ последней, префикс этого
          решения оптимален для состояния-предшественника.
    \item \textbf{Ацикличность}: каждый переход увеличивает хотя бы
          одну координату вектора выпуска, поэтому граф состояний
          ацикличен и динамическое программирование по слоям
          корректно.
    \item \textbf{Полнота}: все максимальные карты сгенерированы;
          доминируемые удалены только в режиме <<не меньше>>, где это
          не влияет на оптимум.
\end{enumerate}





\clearpage

\section{Практическая часть}

Исходные данные:


Лист материала: $3 \times 5$ (15 клеток). Типы фигур и потребность в них
приведены в таблице~\ref{tab:figures}.

\begin{table}[h!]
\centering
\caption{Типы фигур и потребность}
\label{tab:figures}
\begin{tabular}{|c|c|c|c|}
\hline
Фигура & Размер & Площадь & Требуется, шт \\
\hline
A & $3 \times 2$ & 6 & 21 \\
\hline
B & $2 \times 2$ & 4 & 15 \\
\hline
C & $1 \times 3$ & 3 & 22 \\
\hline
D & $1 \times 2$ & 2 & 31 \\
\hline
\end{tabular}
\end{table}

Режим задачи: <<не меньше>> (допускается перевыполнение плана).


В результате полного перебора максимальных размещений было получено
21 карта раскроя. Ниже приведены их описания в виде векторов
$a_j = (a_{jA}, a_{jB}, a_{jC}, a_{jD})$ (таблица~\ref{tab:cards}).

\begin{table}[h!]
\centering
\caption{Карты раскроя и количество фигур на каждой}
\label{tab:cards}
\small
\begin{tabular}{|c|c|c|c|c||c|c|c|c|c|}
\hline
Карта & A & B & C & D & Карта & A & B & C & D \\
\hline
1  & 0 & 0 & 1 & 6 & 12 & 0 & 1 & 0 & 5 \\
\hline
2  & 2 & 0 & 1 & 0 & 13 & 0 & 2 & 0 & 3 \\
\hline
3  & 1 & 1 & 1 & 1 & 14 & 0 & 2 & 2 & 0 \\
\hline
4  & 1 & 0 & 3 & 0 & 15 & 1 & 0 & 0 & 4 \\
\hline
5  & 1 & 0 & 1 & 3 & 16 & 0 & 0 & 4 & 1 \\
\hline
6  & 0 & 2 & 1 & 2 & 17 & 1 & 0 & 2 & 1 \\
\hline
7  & 0 & 1 & 3 & 1 & 18 & 1 & 1 & 0 & 2 \\
\hline
8  & 0 & 1 & 1 & 4 & 19 & 0 & 0 & 2 & 4 \\
\hline
9  & 0 & 0 & 5 & 0 & 20 & 2 & 0 & 0 & 1 \\
\hline
10 & 0 & 0 & 3 & 3 & 21 & 0 & 0 & 0 & 7 \\
\hline
11 & 0 & 1 & 2 & 2 &    &   &   &   &   \\
\hline
\end{tabular}
\end{table}

Всего 21 карта. Каждая карта максимальна: ни одна из фигур A--D ни в одном
из двух допустимых поворотов не может быть добавлена на лист.


После применения метода динамического программирования по вектору выпуска
$(A, B, C, D) = (21, 15, 22, 31)$ получено оптимальное решение:
минимум листов равен $\mathbf{21}$.

Оптимальный план раскроя приведён в таблице~\ref{tab:plan}.

\begin{table}[h!]
\centering
\caption{Оптимальный план раскроя}
\label{tab:plan}
\begin{tabular}{|c|c|c|l|}
\hline
Карта & Количество листов & Вектор выпуска & Состав \\
\hline
$x_1$    & 3  & $(0, 0, 1, 6)$ & C=1, D=6 \\
\hline
$x_2$    & 10 & $(2, 0, 1, 0)$ & A=2, C=1 \\
\hline
$x_3$    & 1  & $(1, 1, 1, 1)$ & A=1, B=1, C=1, D=1 \\
\hline
$x_6$    & 6  & $(0, 2, 1, 2)$ & B=2, C=1, D=2 \\
\hline
$x_{14}$ & 1  & $(0, 2, 2, 0)$ & B=2, C=2 \\
\hline
\end{tabular}
\end{table}

Проверим выполнение плана:
\begin{align*}
A &: 10 \cdot 2 + 1 \cdot 1 = 21 \ge 21, \\
B &: 1 \cdot 1 + 6 \cdot 2 + 1 \cdot 2 = 15 \ge 15, \\
C &: 3 \cdot 1 + 10 \cdot 1 + 1 \cdot 1 + 6 \cdot 1 + 1 \cdot 2 = 22 \ge 22, \\
D &: 3 \cdot 6 + 1 \cdot 1 + 6 \cdot 2 = 31 \ge 31.
\end{align*}

Все ограничения выполнены точно (перевыполнения нет). Суммарное число
использованных листов:
\[
3 + 10 + 1 + 6 + 1 = 21.
\]


В таблице~\ref{tab:result} сопоставлены требуемое и полученное количество
фигур.

\begin{table}[h!]
\centering
\caption{Сравнение плана и факта}
\label{tab:result}
\begin{tabular}{|c|c|c|}
\hline
Фигура & Требуется & Получено \\
\hline
A & 21 & 21 \\
\hline
B & 15 & 15 \\
\hline
C & 22 & 22 \\
\hline
D & 31 & 31 \\
\hline
\end{tabular}
\end{table}



Оптимальный план использует пять карт: №~1, №~2, №~3, №~6 и №~14.
Ниже для каждой карты приведено её графическое представление
(буквами обозначены клетки, занятые соответствующей фигурой;
точка~<<.>> обозначает пустую клетку) и состав.

\clearpage

\begin{figure}[h!]
\centering

% ---------- Карта 1 ----------
\begin{minipage}[t]{0.18\textwidth}
\centering
\textbf{Карта 1}\\[2mm]
{\ttfamily
\begin{tabular}{cccc}
C & D & D & \\
C & D & D & \\
C & D & D & \\
D & D & D & \\
D & D & D & \\
\end{tabular}
}\\[2mm]
\small C=1, D=6\\[1mm]
\small (3 листа)
\end{minipage}\hfill
% ---------- Карта 2 ----------
\begin{minipage}[t]{0.18\textwidth}
\centering
\textbf{Карта 2}\\[2mm]
{\ttfamily
\begin{tabular}{cccc}
A & A & A & \\
A & A & A & \\
A & A & A & \\
A & A & A & \\
C & C & C & \\
\end{tabular}
}\\[2mm]
\small A=2, C=1\\[1mm]
\small (10 листов)
\end{minipage}\hfill
% ---------- Карта 3 ----------
\begin{minipage}[t]{0.18\textwidth}
\centering
\textbf{Карта 3}\\[2mm]
{\ttfamily
\begin{tabular}{cccc}
A & A & A & \\
A & A & A & \\
B & B & C & \\
B & B & C & \\
D & D & C & \\
\end{tabular}
}\\[2mm]
\small A=1, B=1, C=1, D=1\\[1mm]
\small (1 лист)
\end{minipage}\hfill
% ---------- Карта 6 ----------
\begin{minipage}[t]{0.18\textwidth}
\centering
\textbf{Карта 6}\\[2mm]
{\ttfamily
\begin{tabular}{cccc}
B & B & C & \\
B & B & C & \\
B & B & C & \\
B & B & D & \\
D & D & D & \\
\end{tabular}
}\\[2mm]
\small B=2, C=1, D=2\\[1mm]
\small (6 листов)
\end{minipage}\hfill
% ---------- Карта 14 ----------
\begin{minipage}[t]{0.18\textwidth}
\centering
\textbf{Карта 14}\\[2mm]
{\ttfamily
\begin{tabular}{cccc}
B & B & C & \\
B & B & C & \\
B & B & C & \\
B & B & . & \\
C & C & C & \\
\end{tabular}
}\\[2mm]
\small B=2, C=2\\[1mm]
\small (1 лист)
\end{minipage}

\caption{Карты раскроя, вошедшие в оптимальный план}
\label{fig:selected-cards}
\end{figure}
Здесь фигура A занимает прямоугольник $3 \times 2$, B -- $2 \times 2$,
C -- вертикальную полосу $1 \times 3$, D -- $1 \times 2$. Карта
максимальна -- ни одна из фигур не может быть добавлена.

\clearpage

\section{Заключение}

В ходе выполнения лабораторной работы была решена задача оптимального
раскроя листа $3 \times 5$ фигурами четырёх типов A, B, C, D с заданной
потребностью. Были выполнены следующие этапы:

\begin{enumerate}
    \item Сгенерированы все 21 максимальная карта раскроя с учётом
          допустимых поворотов фигур. Использован рекурсивный перебор с
          мемоизацией состояний листа.
    \item Применено отсеивание доминируемых карт в режиме.
    \item Поставлена и решена задача целочисленного программирования
          минимизации числа листов методом динамического программирования
          по четырёхмерному вектору выпуска.
    \item Получен оптимальный план: 21 лист, из них 3 листа по карте~1,
          10 листов по карте~2, 1 лист по карте~3, 6 листов по карте~6
          и 1 лист по карте~14.
\end{enumerate}

Все ограничения по выпуску выполнены точно, минимальное число листов
равно 21. Полученные результаты демонстрируют эффективность сочетания
полного перебора схем раскроя с методом динамического программирования
для задач дискретной оптимизации. Цель лабораторной работы достигнута.

\clearpage

\section{Список литературы}

\begin{enumerate}
    \item Канторович Л.В., Залгаллер В.А. Рациональный раскрой
          промышленных материалов. -- Новосибирск: Наука, 1971. -- 300 с.
    \item Gilmore P.C., Gomory R.E. A linear programming approach to the
          cutting-stock problem // Operations Research. -- 1961. --
          Vol.~9, No.~6. -- P.~849--859.
    \item Беллман Р. Динамическое программирование. -- М.: Изд-во
          иностранной литературы, 1960. -- 400 с.
    \item Кормен Т., Лейзерсон Ч., Ривест Р., Штайн К. Алгоритмы:
          построение и анализ. -- М.: Вильямс, 2013. -- 1328 с.
    \item Лубенец Ю.В. Теория принятия решений: Учеб. пособие. --
          Липецк: Изд-во ЛГТУ, 2011. -- 80 с.
\end{enumerate}
